\chapter{Decomposição de Séries II}
\label{cap:filtros2}
% ── índice ──────────────────────────────────────────────────────
\index{STL|textbf}
\index{stl@\texttt{stl()}|textbf}
\index{Christiano-Fitzgerald, filtro de|textbf}
\index{christiano_fitzgerald@\texttt{christiano\_fitzgerald()}|textbf}
\index{loess}
\index{filtro passa-banda}
\index{decomposição sazonal}

% ─────────────────────────────────────────────────────────────────────────────
\section{STL e Christiano-Fitzgerald}

O cap.~\ref{cap:filtros} apresentou HP, BK e Holt-Winters. Este capítulo
cobre dois métodos adicionais:

\begin{description}
  \item[STL] \emph{Seasonal-Trend decomposition via Loess} (Cleveland et al.~1990):
    decompõe $y_t = T_t + S_t + R_t$ iterativamente via regressão local (loess).
    Robusto a outliers; permite sazonalidade variante no tempo.

  \item[Christiano-Fitzgerald (CF)] Filtro passa-banda assimétrico (2003):
    extrai frequências em $[p_L, p_H]$ usando pesos de uma janela
    não-simétrica que se adapta à posição na série. Ao contrário do BK, não
    perde pontos de borda — o preço é que os pesos mudam a cada observação.
\end{description}

\subsection*{STL: decomposição iterativa}

O algoritmo STL alterna entre:
\begin{enumerate}
  \item \emph{Seasonal loop}: suaviza cada posição sazonal com loess
  \item \emph{Trend loop}: suaviza a série dessazonalizada com loess
\end{enumerate}
A convergência produz $T_t$ (tendência), $S_t$ (sazonalidade) e $R_t$ (resíduo)
com propriedades de robustez a outliers quando \texttt{robust=true}.

\subsection*{Filtro Christiano-Fitzgerald}

O CF aplica pesos $a_j$ que satisfazem $\sum_j a_j e^{i\omega j} \approx 1$
para $\omega \in [\omega_L, \omega_H]$ e $\approx 0$ fora da banda.
Para $t$ nas bordas, os pesos se ajustam para não usar observações futuras.

% ─────────────────────────────────────────────────────────────────────────────
\section{Verificação por DGP}

DGP: série mensal com tendência linear, ciclo anual e sazonalidade semestral.
$T = 120$ meses.

\begin{lstlisting}[caption={STL e Christiano-Fitzgerald}]
seed(42)
generate df trend  = 0.05 * t
generate df cycle  = 0.5 * sin(t * 3.14159 / 12)
generate df season = 0.3 * sin(t * 3.14159 / 6)
generate df y      = 10.0 + trend + cycle + season + noise

stl(df, y, period=12)
display df

christiano_fitzgerald(df, y, low=6, high=32)
display df
\end{lstlisting}

\subsection*{STL}

\begin{lstlisting}[language={}, basicstyle=\ttfamily\small, frame=none,
  numbers=none, backgroundcolor=\color{codbg}]
Seasonal Decomposition (STL)
  Observations: 120
  Trend mean:   13.0243
  Seasonal std:  0.1943
  Residual std:  0.1868
\end{lstlisting}

O STL adiciona colunas \texttt{y\_trend}, \texttt{y\_seasonal} e \texttt{y\_resid}
ao DataFrame. Residual std $< 0{,}2$ — o ruído verdadeiro ($\sigma = 0{,}3$) é
bem separado dos demais componentes.

\subsection*{Christiano-Fitzgerald}

\begin{lstlisting}[language={}, basicstyle=\ttfamily\small, frame=none,
  numbers=none, backgroundcolor=\color{codbg}]
cffilter: períodos [6,32]  →  y_cycle adicionada a df
  t=1:  y_cycle = 1.04
  t=6:  y_cycle = 1.19    (ciclo anual dominante)
  t=12: y_cycle = 0.87
\end{lstlisting}

O CF retém variações periódicas entre 6 e 32 meses (1{,}5 a 8 anos — ciclos
de negócios). A vantagem sobre o BK é que todas as 120 observações são mantidas
(sem trimagem de pontas).

% ─────────────────────────────────────────────────────────────────────────────
\section{Sintaxe}

\begin{lstlisting}
stl(df, var, period=12)                        // adiciona trend, seasonal, resid
christiano_fitzgerald(df, var, low=6, high=32) // adiciona var_cycle
\end{lstlisting}

\begin{center}
\begin{tabular}{llll}
\toprule
\textbf{Filtro} & \textbf{Tipo} & \textbf{Perda de borda} & \textbf{Sazonalidade} \\
\midrule
HP              & passa-baixo    & Sim (ponta)         & Não \\
BK              & passa-banda    & Sim ($k$ cada lado) & Não \\
CF              & passa-banda    & Não                 & Não \\
STL             & decomposição   & Não                 & Sim \\
\bottomrule
\end{tabular}
\end{center}

\begin{nota}
  A função \texttt{hamilton} no Hayashi refere-se ao modelo de Markov
  Switching de Hamilton (1989), \emph{não} ao filtro de Hamilton (2018)
  (cap.~\ref{cap:regime}). O filtro Hamilton 2018 — que propõe substituir
  o HP por uma regressão OLS com 4 defasagens e horizonte $h = 8$ —
  não está implementado separadamente.
\end{nota}
